\documentclass[instructions.tex]{subfiles}
\begin{document}
 
\chapter{Du paradis.}
 
Qu’est-ce que le paradis? Que faut-il faire pour le mériter?
 
\primo Le paradis est le séjour de la magnificence et de la gloire de Dieu, la demeure qu’il a préparée à ses anges, et à ceux qui ont vécu dans sa crainte, et qui meurent dans son amour. C’est le lieu le plus beau, le plus riche, le plus délicieux, le plus noble, le plus vaste; en un mot le chef-d’œuvre de la puissance du Très-Haut.
 
Nous pouvons dire ce qui n’est pas dans le paradis; nous pouvons dire quelque chose de ce qui y est; mais tout ce que nous en disons, c’est comme si nous n’en disions rien. Il n’y a rien dans le ciel de tout ce qu’on voit ici-bas: il n’y a ni faim, ni soif, ni froid, ni chaud, ni adversités, ni maladies, ni lassitudes, ni ennuis: les plaisirs, les honneurs qui y sont, l’éclat et les richesses qu’on y voit, n’ont rien de semblable à tout ce qu’on voit sur la terre.
 
Le paradis surpasse tellement nos pensées, que saint Paul, à qui Jésus-Christ en fit voir quelques rayons, ne nous en peut dire autre chose, sinon que l’œil n’a rien vu, l’oreille n’a point entendu, l’esprit de l’homme n’a point conçu ce que Dieu a préparé à ceux qui l’aiment\footnote{Oculus non vidit, nec auris audivit, nec in cor hominis ascendit, quæ præparavit Deus iis qui diligunt illum. 1. Cor. 2. 9.}. Nos corps y seront plus brillans que le soleil; immortels et impassibles, il ne pourront ni souffrir, ni vieillir, ni mourir; ils seront tous doués d’une santé, d’une beauté, d’une agilité incorruptibles; ils seront inondés de délices, et goûteront des plaisirs si purs et si saints, que notre imagination ne peut rien y comprendre ici-bas.
 
Mais le bonheur souverain de l’ame dans le ciel, c’est de voir et de posséder Dieu. Ce bonheur, ne fût-il que pour un moment, est quelque chose de si élevé au-dessus des mérites d’une créature, que quand vous auriez souffert les tourmens des martyrs pour le mériter, ce serait une présomption d’y prétendre sans la promesse que Dieu nous en a faite. Être dans le ciel, et y voir Dieu pendant un quart d’heure, est un plus grand bonheur que de posséder tous les royaumes de la terre pendant un siècle. Que sera-ce de voir et de posséder Dieu pour toujours?
 
II n’est aucun bien dans le ciel qui n’y soit dans un souverain degré: la gloire, les plaisirs, la santé, tout ce qui peut rendre heureux y est infini. Le cœur y est toujours dans une douce extase, dans le ravissement; toujours inondé de ces torrens de délices et de joie qui découlent de la vue de Dieu, sans interruption et sans dégoût.
 
Si un damné était assuré qu’après cent ans de tourmens, il jouirait pendant une heure des délices du ciel et de la possession de Dieu, il serait consolé. Nous espérons de posséder Dieu et de jouir de sa félicité, non pas pour quelques momens, mais pour toujours. Loin de travailler pour acquérir ce bonheur, nous n’y pensons même pas. Où sont notre raison et notre foi?
 
\secundo Si nous ne pouvons comprendre ce que c’est que le paradis, nous pouvons du moins nous disposer à le mériter. Je dis nous disposer, car personne ne peut le mériter par soi-même en rigueur. Nous nous y disposons par nos bonnes œuvres, et nous l’obtenons par les mérites du Sauveur.
 
La porte en est étroite, dit Jésus-Christ, faites vos efforts pour y entrer. Quelque grands que doivent être ces efforts, ils sont peu de chose en comparaison d’un si grand bien. Nous ne devons pas nous affliger de la peine qu’il faut prendre pour l’obtenir: les bienheureux regardent comme un rien tout ce qu’ils ont souffert pour acquérir un bien si désirable. Ce ne serait pas trop, dit saint Augustin, d’une éternité de travaux pour s’assurer une éternité de bonheur\footnote{Æterna requies æterno labore rectè emitur. Aug.}.
 
Le ciel est une récompense, il faut donc travailler pour la gagner; c’est une couronne, il faut combattre pour l’obtenir; c’est une conquête, il faut la violence pour la mériter. Ceux qui savent se vaincre eux-mêmes, sont les héros qui l’emportent\footnote{Violenti rapiunt illud. Matth. 11. 12.}.
 
Le ciel est un séjour de sainteté; les portes de cette cité sainte ne sont ouvertes qu’à l’innocence ou à la pénitence. Si vous n’êtes occupé que des vanités, des biens et des plaisirs de cette vie, vous n’êtes ni innocent, ni pénitent; vous êtes donc indigne d’entrer dans le ciel; vous êtes même indigne de porter le nom de chrétien. Vous vivez comme si vous ne croyiez point; car peut-on dire que vous croyez le paradis, lorsque vous ne travaillez pas à le mériter? Un homme croirait-il qu’il y a un trésor dans ses fonds, s’il ne voulait pas seulement se donner la peine d’y faire creuser pour le trouver?
 
Si vous étiez assuré qu’un puissant roi enrichit ceux qui vont habiter dans ses terres, vous quitteriez tout pour aller vous y établir. Voilà ce que vous feriez pour aller vous enrichir dans un royaume où vous n’auriez tout au plus que vingt ou trente ans à vivre, tandis que vous ne faites rien pour vous établir pour toujours dans le ciel. En vérité, avez-vous la foi? La marque qu’un chrétien croit et espère le ciel, c’est quand il méprise les choses de la terre, et tout ce qui se passe dans cette vie, dit saint Grégoire. Que sert-il de croire qu’il y a un paradis, si l’on ne fait rien pour le gagner, et si l’on fait tout pour le perdre?
 
\section{Résolutions}
 
\primo Tous les matins je me dirai, en m’habillant: Voici une journée que Dieu m’accorde pour travailler à mériter la place qu’il m’a marquée dans le paradis: je dois donc accomplir tous mes devoirs, et me préserver du péché.
 
\secundo Je regarderai le paradis comme ma patrie, vers laquelle je dois tendre de toutes mes forces.
 
\prayer{Mon Dieu, aidez-moi du secours de votre grâce, afin que je marche avec tant de précaution, pendant mon court pèlerinage sur la terre, que j’aie le bonheur d’aller un jour régner avec vous dans le séjour immortel de la gloire.}
 
\end{document}
